\documentclass[printMode=true, declarePage=false]{ecnuthesis}
% 
%
% ECNUThesis-Undergraduate: 华东师范大学本科生毕业论文模板
% 2024-04-16 v0.2
%
% 本模板源代码托管在
% https://github.com/Koyamin/ECNUThesis-Undergraduate
% 最新的 Release 版本将发布在该仓库
% 
\ecnuSetup{
  info = {
    title = {基于微信小程序的线下课程考勤系统的设计与实现},
    % 中文标题
    %
    titleEN = {Design and Implementation of Offline Course Attendance System Based on WeChat Mini Program},
    % 英文标题
    %
    author = {张三},
    % 作者姓名
    %
    studentID = {10000000000},
    % 作者学号
    %
    department = {理工学院},
    % 学院名称
    %
    major = {计算机科学与技术},
    % 专业名称
    %
    supervisor = {李四},
    % 指导教师姓名
    %
    academicTitle = {教授},
    % 指导教师职称
    %
    year  = 2077,
    % 论文完成年份
    %
    month = 3,
    % 论文完成月份
    %
    graduationYear = 2077,
    % 内封面毕业届别
    % 说明:
    %   若 graduationYear 字段为空，则内封面毕业届别为 year
    %
    keywords = {XML, XSL, 数据批处理, “双赢”},
    % 中文关键词
    % 请使用英文逗号 "," 以分隔
    %
    keywordsEN = {XML, XSL, batch, ‘Win-Win’},
    % 英文关键词
    % 请使用英文逗号 "," 以分隔
    %
  },
  style = {
    numbering = arabic,
    % 章节编号样式
    % 可用选项:
    %   numbering = arabic|alpha|chinese
    % 说明:
    %   arabic    使用数字进行编号 (即理科要求)
    %   alpha     使用字母进行编号 (即外文要求)
    %   chinese   使用汉字进行编号 (即文科要求)
    %   (默认选项为 arabic )
    %
    font = xits,
    % 西文字体选择
    % 可用选项:
    %   font = lm|times|xits
    % 说明:
    %   lm        使用 TeX 自带的 Latin Modern 字体
    %   lm        使用 TeX 自带的 XITS 系列字体
    %   times     使用 Times 风格的西文字体
    %   (默认选项为 xits )
    %
    fontCJK = fandol,
    % 中文字体选择
    % 可用选项:
    %   fontCJK = fandol|windows|mac
    % 说明:
    %   fandol    使用 TeX 自带的 fandol 字体
    %   windows   使用 Windows 系统内的字体 (中易)
    %   mac       使用 MacOS 系统内的字体
    %   (默认选项为 fandol )
    %
    logoResource = {./ecnu-vi-inner-cover-logo.pdf},
    % 封面插图数据源
    % 模版已自带, 位于 ./source/inner-cover(contains_font).eps
    % 默认值为空
    %
    % declarePageResource = {./source/declaration.pdf},
    % 扫描版声明页 PDF 文件
    % 若该值为空则生成模版预定义的声明页；否则将插入指定路径所对应的 PDF 文件
    % 默认值为空
    %
    bibResource = {./bib/thesis-ref.bib},
    % 参考文献数据源
    % 由于使用的是 biber + biblatex , 所以必须明确给出 .bib 后缀名
    %
  }
}

\usepackage{enumitem}
\usepackage{mwe}
\usepackage{zhlipsum}

\begin{document}

% 目录
\tableofcontents

% 声明前置部分开始, 页码设置为大写罗马数字
\frontmatter

\begin{abstract}
  本文通过一个实际的对日软件外包案件的设计和实现，经历了整个软件开发的过程，包括系统分析、概要设计、详细设计、编码、测试，为某制药企业开发了一个 B2B 的电子商务系统。本系统主要是以该制药企业为购买方，发布企业所需要的货物清单，以本系统为平台，各个供应商进行竞标，由购买方选择购买供应商的货物并下订单发货。由于购买方发布货物需要对大量数据进行操作，因此制作了数据批处理程序，来实现大量数据的导入和导出。系统的在线部分运用了XML和XSL技术，体现了画面实现模版化的优势，使得更有效方便的实现画面重载。该系统的开发过程中，本人主要负责制造在线部分的登陆、Home、估价请求履历检索、取消订单一览、用户信息作成、通知信息作成6个模块以及数据批处理部分的供应商 Master 导入程序。

  使用该系统，制药企业可以在众多供应商中选择最价廉物美的原材料，这样大大降低其成本，提高了企业利润。同时，供应商之间也有了相互竞争，可以促进生产，达到“双赢”的效果。
  本文最后说明了对日软件开发过程与当今我国软件开发过程的区别，并对我国今后软件事业做了期望和展望。
\end{abstract}

\begin{abstractEN}
  This paper concerns with design and achievement of a software development project for Japanese company, experence the whole process of software development including system analysis, general design, detail design, coding, testing to develop a B2B E-business system for some Medicine Manufacture Enterprise. In this system, the Medicine Manufacture Enterprise as the buyers publish products list.The suppliers use this system as plat to bit products.Then the buyers select suppliers they need and send order to buy products.This system also has some batches to realize huge data importing and exporting.This system's webs use XML and XSL technologies to realize reloading of web templates. During the development of the system, my duty is to develop six webs and one batch.

  Using this system, Medicine Manufacture Enterprise can select the best and the cheapest products in several suppliers to descreat cost. At the same time, suppliers will bit each other to inprove production and reach ‘Win-Win’.
  
  At the end of this paper, also explain the deference between Janpanese software development processing and Chinese, and give the deep expectance.
\end{abstractEN}

% 声明正文部分开始, 页码设置为阿拉伯数字
\mainmatter

\chapter{绪论}

\section{背景}

\subsection{浅谈中国软件}

众所周知，信息产业是……

\subsection{项目背景}

如今，随着互联网的日益流行，……

\subsection{关于Microsoft Commerce Server 2002}

\subsubsection{简介}

\subsubsection{技术特点}

Microsoft Commerce Server 2002 拥有以下技术特点\footnote{技术特点是指…}。

\paragraph{国际性站点支持}

Microsoft Commerce Server 2002 对国际性站点有着良好的支持……

\paragraph{更强大的B2B}

与其余同类型技术相比，Microsoft Commerce Server 2002 对B2B的支持……


\subsection{关于Commerce Brains}

\subsubsection{项目介绍}

\clearpage

所谓 Commerce Brains，就是 B to B 专用的、以 web 为基础的应用程序包\cite{book:begin-c-sharp-xml}。
Commerce Brains 是由业务组件\cite{book:html-and-xml-web}和 Enterprise Manager \cite{book:ast-web-programming}构成的，如图 \ref{fig:component-deployment} 
组件分布图所示。

\begin{figure}[htb]
  \centering
  \includegraphics[width=.5\textwidth]{example-image}
  \bicaption{组件分布图}{Component Deploment}\label{fig:component-deployment}
\end{figure}

\subsubsection{产品特点}

\zhlipsum[1]

\subsection{相关工作}

\zhlipsum[2]

\subsection{论文组织}

\zhlipsum[3]


\chapter{系统分析与设计}

\section{需求分析}

\zhlipsum[4]

\section{系统概要}

\zhlipsum[5]

\section{系统流程}

\zhlipsum[6]


\chapter{概要设计}

\section{业务流程}

\zhlipsum[7]

\section{功能模块介绍}

如图 \ref{fig:system-use-case} 系统用例图所示，系统的业务流程分为两部分，购买 WF 和 In 购买。

\begin{figure}[htb]
  \centering
  \includegraphics[width=.5\textwidth]{example-image}
  \bicaption{系统用例图}{Component Deploment}\label{fig:system-use-case}
\end{figure}


\chapter{详细设计}

\section{数据库的设计}

\zhlipsum[8]

\section{详细设计书}

\zhlipsum[9]

\section{界面设计}

\zhlipsum[10]


\chapter{系统实现}

\section{界面}

\zhlipsum[11]


\chapter{关键算法与数据结构}

\zhlipsum[12]


\chapter{系统测试}

\section{测试计划}

\subsection{单元测试}

单元测试主要是对画面显示的测试，主要测试点如表7-1所示：

单元测试主要是对画面显示的测试，主要测试点如表 \ref{table: unit-test-case} 所示。

\begin{table}[htb]
\centering
\bicaption{单体测试表}{Unit Test Case}\label{table: unit-test-case}
\begin{tabular}{|l|l|}
  \hline
  输入系 & \begin{tabular}[c]{@{}l@{}}输入正常值\\ （最小值\textless{}值\textless{}最大值）\end{tabular} \\ \hline
      & \begin{tabular}[c]{@{}l@{}}输入系统异常值\\ （值=最小值-1，值=最大值+1）\end{tabular}  \\ \hline
      & \begin{tabular}[c]{@{}l@{}}输入系统正常值（边际测试）\\ （值=最小值，值=最大值）\end{tabular} \\ \hline
      & 未输入值 \\ \hline
  显示系 & 标题（Title）显示                                                                     \\ \hline
      & 说明文显示                                                                           \\ \hline
      & 显示相应业务的左菜单  \\ \hline
      & 显示Text项目（不换行） \\ \hline
      & 按需要显示排序按钮 \\ \hline
      & 添加附件参照  \\ \hline
  操作系 & IE返回按钮按下后画面迁移                                                                   \\ \hline
      & IE $\times$按钮按下后，系统终了                                                           \\ \hline
      & 下一画面迁移按钮按下后的迁移                                                                  \\ \hline
      & 上一画面迁移按钮按下后的迁移                                                                  \\ \hline
\end{tabular}
\end{table}

\subsection{结合测试}

\zhlipsum[12]

\subsection{模拟测试}

\zhlipsum[13]

\subsection{随机测试}

\zhlipsum[14]


\chapter{测试用例与测试报告}

\zhlipsum[15]


\chapter{总结与展望}

\zhlipsum[16]

% 声明后置部分开始
\backmatter

% 致谢
\begin{acknowledgement}
  感谢天，感谢地，感谢阳光照耀了大地。
\end{acknowledgement}

% 打印参考文献
\PrintReference

% 附录
\begin{appendix}
  \begin{enumerate}
    \item 部分函数代码
    \begin{enumerate}[label=（\arabic*）]
      \item LoadMessageFile函数（过程），用于读取信息XML文件，取得相应节点的内容，并通过引用参数返回值。
    \end{enumerate}
  \end{enumerate}
\end{appendix}

\end{document}